%%%%%%%%%%%%%%%%%%%%%%%%%%%%%%%%%%%%%%%%%%%%%%%%%%%%%%%%%
%Modelos de listas que se pueden usar, solo es necesario copiar y pegar el modelo y cambiar el texto del interior, añadiendo o quitando puntos en función de los que sea necesarios. Los puntos se agregan con el comando "\item", seguido del texto que queramos añadir
%%%%%%%%%%%%%%%%%%%%%%%%%%%%%%%%%%%%%%%%%%%%%%%%%%%%%%%%%

\begin{itemize}
    \item Esta es una lista con puntos, es decir, no numerada.
    \item Podemos colocar tantos puntos como queramos.
    \item El tamaño de los textos es indiferente, ya que se adaptan al tamaño del pdf sin ningún problema.
\end{itemize}

\begin{enumerate}
    \item Esta es una lista numerada.
    \item Al igual que ocurría con el anterior se pueden poner textos de cualquier longitud.
    \item También podemos poner tantos números como queramos.
\end{enumerate}